\begin{theoremEnd}[normal,category=combinations]{corollary}[Signed combinations and their reductions]
\CombinedResultCheck{cor:combinations}{corollary}\label{cor:combinations}
Under the finite Gaussian model in \cref{sec:finite-model}, let
$A\subseteq\{0,1,2\}$ be nonempty and fix signs
$\varepsilon_r\in\{-1,1\}$ for $r\in A$.
For $\theta(x)=\omega_1(x-c+L)$, $w(x)=\cos\theta(x)$, and
$J=[w(x_+),w(x_-)]\subset(-1,1)$, define, for a coefficient array
$\boldsymbol b=(b_0,\ldots,b_m)^T$,
\[
 \begin{aligned}
 H_{0,\boldsymbol b}(w)
   &=\sqrt L\,b_0+\sqrt{1-w^2}\sum_{j=1}^m b_jU_{j-1}(w),\\
 H_{1,\boldsymbol b}(w)&=\sum_{j=1}^m jb_jT_j(w),&
 H_{2,\boldsymbol b}(w)&=-\sum_{j=1}^m j^2b_jU_{j-1}(w).
 \end{aligned}
\]
Suppose a fixed $\boldsymbol d=(d_0,\ldots,d_m)^T$ satisfies
$\min_{w\in J}\varepsilon_rH_{r,\boldsymbol d}(w)>0$ for every
$r\in A$. Set $h=d_0+\sum_{j=1}^m d_j\phi_{j,L}$ and choose
$j_*\in\{0,\ldots,m\}$ with $d_{j_*}\ne0$. The coordinates
\[
 \gamma=\frac{a_{j_*}}{d_{j_*}},\qquad
 g_j=a_j-d_j\gamma\quad(j\ne j_*),\qquad g_{j_*}=0
\]
give $\boldsymbol g=(g_j)_{j\ne j_*}\in\R^m$ and
$f_{\boldsymbol g}=g_0+\sum_{j=1}^m g_j\phi_{j,L}$.
Then $f_{\boldsymbol a}=f_{\boldsymbol g}+\gamma h$, with inverse
$a_j=g_j+d_j\gamma$ for $0\leq j\leq m$.
In expressions involving $H_{r,\boldsymbol g}$, the zero coordinate
$g_{j_*}$ is included in the coefficient array.

With $\eta_r$ from \cref{eq:joint-threshold}, define
\[
 \eta_r(x;\boldsymbol g)
   =-\frac{H_{r,\boldsymbol g}(w(x))}{H_{r,\boldsymbol d}(w(x))},
 \qquad M_r(\boldsymbol g)=\max_{x\in I}\eta_r(x;\boldsymbol g),
 \qquad \eta_r^\pm(\boldsymbol g)=\eta_r(x_\pm;\boldsymbol g).
\]
For adjacent selected orders $r,r+1\in A$, let
$E_r(\boldsymbol g)=\eta_r^-(\boldsymbol g)$ when
$\varepsilon_r=\varepsilon_{r+1}$ and
$E_r(\boldsymbol g)=\eta_r^+(\boldsymbol g)$ otherwise.
Suppressing the argument $\boldsymbol g$ on the right, the threshold is
\[
 \tau_A^\varepsilon(\boldsymbol g)=
 \begin{cases}
 M_r, & A=\{r\},\\
 \max\{M_1,E_0\}, & A=\{0,1\},\\
 \max\{M_2,M_0\}, & A=\{0,2\},\ \varepsilon_0=\varepsilon_2,\\
 \max\{M_2,\eta_0^-,\eta_0^+\},
     & A=\{0,2\},\ \varepsilon_0\ne\varepsilon_2,\\
 \max\{M_2,E_1\}, & A=\{1,2\},\\
 \max\{M_2,E_1,E_0\}, & A=\{0,1,2\}.
 \end{cases}
\]
\cref{prop:joint} applies with scalar coefficient
$\gamma=a_{j_*}/d_{j_*}$. Thus
$f_{\boldsymbol a}\in\mathcal C_A^\varepsilon$ exactly when
$\gamma\geq\tau_A^\varepsilon(\boldsymbol g)$.

Each maximum $M_r$ is attained at a physical endpoint or a location
whose cosine coordinate is an interior root in $J$ of
\[
 H_{r,\boldsymbol d}(w)\frac{\mathrm{d}H_{r,\boldsymbol g}}{\mathrm{d}w}(w)
 -H_{r,\boldsymbol g}(w)\frac{\mathrm{d}H_{r,\boldsymbol d}}{\mathrm{d}w}(w)=0.
\]
For $r=1,2$, the left side is a polynomial.
All root multiplicities and reduced degrees are included.
If the left side vanishes identically, the quotient is constant and
either endpoint suffices.

These formulas cover the 26 nonempty combinations without opposing
signs, with positive ordinary posterior probability whenever the
stated direction condition holds for the selected combination.
The remaining 37 selections contain opposing signs. At orders zero,
one, and two these force, respectively, a zero, constant, or affine
function, subject to any further restrictions. In the finite
expansion, affine curves are constant, so all 37 events have
ordinary posterior probability zero.
\end{theoremEnd}

\begin{proofEnd}
\CombinedProofCheck{cor:combinations}
Fix the selected orders, signs, and direction coefficients satisfying
\cref{cor:combinations}. The inverse map follows by substitution, so
the coordinates preserve the full coefficient space. The Chebyshev identities in
\cref{def:chebyshev} give
\[
 \begin{aligned}
 f_{\boldsymbol b}(x)&=\frac{H_{0,\boldsymbol b}(w(x))}{\sqrt L},\\
 \frac{\mathrm{d}f_{\boldsymbol b}}{\mathrm{d}x}(x)
    &=\frac{\omega_1}{\sqrt L}H_{1,\boldsymbol b}(w(x)),\\
 \frac{\mathrm{d}^2f_{\boldsymbol b}}{\mathrm{d}x^2}(x)
    &=\frac{\omega_1^2\sin\theta(x)}{\sqrt L}
       H_{2,\boldsymbol b}(w(x)).
 \end{aligned}
\]
Since $\sin\theta$ is positive on compact $I$, the assumed conditions
on $H_{r,\boldsymbol d}$ give the uniform signed derivative bounds
required by \cref{prop:joint}. The positive factors cancel from the
ratios, yielding the displayed expression for $\eta_r$ and the
conditioned posterior.

A signed slope restriction makes the minimum of
$\varepsilon_0 f_{\boldsymbol a}$ occur at $x_-$ when
$\varepsilon_0=\varepsilon_1$, and at $x_+$ otherwise.
This gives $E_0$. Applying the same argument to the signed slope
$\varepsilon_1\frac{\mathrm{d}f_{\boldsymbol a}}{\mathrm{d}x}$
under the curvature restriction gives $E_1$. When
$\varepsilon_0\ne\varepsilon_2$, the signed curve
$\varepsilon_0 f_{\boldsymbol a}$ is concave, so its nonnegativity
is equivalent to nonnegative endpoint values by the chord inequality.
When the signs agree, the signed curve is convex and its minimum may
be interior, so the value restriction retains $M_0$. Together with
the continuum slope and curvature bounds $M_1$ and $M_2$, these
equivalences give every displayed threshold, because each pointwise
restriction is $\gamma\geq\eta_r(x;\boldsymbol g)$.

The cosine coordinate is strictly decreasing on $I$, and the signed
denominators stay positive on compact $J$. The extreme value theorem,
Fermat's theorem, and the quotient rule therefore give the stated
endpoint and stationary candidates. The functions $H_{r,\boldsymbol b}$
are analytic on $(-1,1)$, so a nonzero stationary expression has only
finitely many roots on $J$. If it vanishes identically, the mean value
theorem gives constancy of the quotient.

Choosing no restriction or one of two signs at each order gives
$3^3-1=26$ nonempty combinations: six singles, twelve pairs, and eight
triples. Of the $2^6-1=63$ nonempty selections from the six constraints,
the other 37 contain opposing signs, which force the corresponding
derivative to vanish identically. This gives zero, a constant, or an
affine function. Analyticity extends any affine identity on $I$ to
the padded interval, whose two endpoint values both equal $a_0$, so
the affine curve is constant. Linear independence of the sine basis
gives $a_1=\cdots=a_m=0$, with $a_0=0$ as well when both value signs
are imposed. These proper linear subspaces have zero probability
under the nondegenerate Gaussian posterior.
\end{proofEnd}
